% Abkürzungen, ggfs. mit korrektem Leerraum 
\newcommand{\bs}{$\backslash$\xspace}
\newcommand{\bspw}{bspw.\xspace} % beispielsweiße
\newcommand{\bzw}{bzw.\xspace} % beziehungsweiße
\newcommand{\ca}{ca.\xspace} % circa
\newcommand{\dahe}{\mbox{d.\,h.}\xspace} % das heißt
\newcommand{\etc}{etc.\xspace} % etcetera (und so weiter (usw))
\newcommand{\eur}[1]{\mbox{#1\,\texteuro}\xspace} % Euro Zeichen
\newcommand{\evtl}{evtl.\xspace} % eventuell
\newcommand{\ggfs}{ggf.\xspace} % gegebenenfalls
\newcommand{\Ggfs}{Ggf.\xspace} % Gegebenenfalls
\newcommand{\gqq}[1]{\glqq{}#1\grqq{}}
\newcommand{\inkl}{inkl.\xspace} % inklusive
\newcommand{\insb}{insb.\xspace} % insbesondere
\newcommand{\uA}{\mbox{u.\,A.}\xspace} % unter Anderem
\newcommand{\usw}{usw.\xspace} % und so weiter
\newcommand{\vgl}{vgl.\xspace} % vergleich
\newcommand{\Vgl}{Vgl.\xspace} % Vergleich
\newcommand{\zB}{\mbox{z.\,B.}\xspace} % zum Beispiel
\newcommand{\ZB}{\mbox{Z.\,B.}\xspace} % zum Beispiel
\newcommand{\og}{\mbox{o.\,g.}\xspace} % oben genannt
\newcommand{\sog}{sog.\xspace} % sogenannt

% Verweise (im Dokument)
% Default  = Name
% K : Kurz = Seite
% V : Voll = Name und Seite
\newcommand{\Abbildung}[1]{Abbildung~\ref{#1}}
\newcommand{\AbbildungV}[1]{Abbildung~\ref{#1}:~\nameref{#1} auf S.~\pageref{#1}}
\newcommand{\Anhang}[1]{Anhang~\ref{#1}:~\nameref{#1} auf S.~\pageref{#1}}
\newcommand{\AnhangK}[1]{Anhang~\ref{#1} auf S.~\pageref{#1}}
\newcommand{\AnhangKapitel}[1]{Kapitel~\ref{#1}:~\nameref{#1} auf S.~\pageref{#1}}
\newcommand{\Tabelle}[1]{Tabelle~\ref{#1}:~\nameref{#1}}
\newcommand{\TabelleK}[1]{Tabelle~\ref{#1} auf S.~\pageref{#1}}
\newcommand{\TabelleV}[1]{Tabelle~\ref{#1}:~\nameref{#1} auf S.~\pageref{#1}}
\newcommand{\Verweis}[1]{Kapitel~\ref{#1}}
\newcommand{\VerweisKL}[1]{Kapitel~\ref{#1} auf Seite~\pageref{#1}}
\newcommand{\VerweisVL}[1]{Kapitel~\ref{#1}:~\nameref{#1} auf Seite~\pageref{#1}}
\newcommand{\VerweisKT}[1]{(Siehe Kapitel~\ref{#1} auf Seite~\pageref{#1})}
\newcommand{\VerweisVT}[1]{(Siehe Kapitel~\ref{#1}:~\nameref{#1} auf Seite~\pageref{#1})}
\newcommand{\VerweisKF}[1]{\footnote{Siehe Kapitel~\ref{#1} auf Seite~\pageref{#1}.}}
\newcommand{\VerweisVF}[1]{\footnote{Siehe Kapitel~\ref{#1}:~\nameref{#1} auf Seite~\pageref{#1}.}}
\newcommand{\Note}[1]{\footnote{#1}}

% Referenzen (außerhalb des Dokument)
\NewDocumentCommand{\Zitat}{ O{} O{} m }{\parencite[#1][#2]{#3}\xspace}
\NewDocumentCommand{\ZitatVgl}{ O{} m }{\parencite[vgl.][#1]{#2}\xspace}

% Styling
\newcommand{\itemd}[2]{\item{\textbf{#1}}\\{#2}} % erzeugt ein Listenelement mit fetter Überschrift
\newcommand{\changefont}[3]{\fontfamily{#1} \fontseries{#2} \fontshape{#3} \selectfont} % einfaches Wechseln der Schrift, z.B.: \changefont{cmss}{sbc}{n}

% fügt Tabellen aus einer TEX-Datei ein
\newcommand{\newTable}[3] % Parameter: caption, label, file
{\begin{table}[htb]
     \centering
     \singlespacing
     \input{table/#3}
     \caption{#1}
     \label{#2}
\end{table}}
\newcommand{\newTabelleAnhang}[1] % Parameter: file
{\begin{center}
     \singlespacing
     \input{table/#1}
\end{center}}


\newcommand{\includegraphicsKeepAspectRatio}[2]{\includegraphics[width=#2\textwidth,height=#2\textheight,keepaspectratio]{#1}}
% Verwendung analog zu \includegraphics
\newlength{\myx} % Variable zum Speichern der Bildbreite
\newlength{\myy} % Variable zum Speichern der Bildhöhe
\newcommand\includegraphicstotab[2][\relax]{%
% Abspeichern der Bildabmessungen
    \settowidth{\myx}{\includegraphics[{#1}]{#2}}%
    \settoheight{\myy}{\includegraphics[{#1}]{#2}}%
% das eigentliche Einfügen
    \parbox[c][1.1\myy][c]{\myx}{%
        \includegraphics[{#1}]{#2}}%
}

% ### Semantische Textverarbeitung ###
\newcommand{\Abk}[1]{\textit{\ac{#1}}} % Abkürzung
\newcommand{\AbkK}[1]{\textit{\acs{#1}}} % Abkürzung kurz (ohne Erstnennung ausgeschrieben)
\newcommand{\Datum}[1]{\DTMdate{#1}} % Datum
\newcommand{\Monat}[1]{% Datum ohne Tag, z.B. \Monat{2020-09}
    \DTMlangsetup{showdayofmonth=false}%
    \DTMdate{#1-01}%
    \DTMlangsetup{showdayofmonth=true}%
}
\newcommand{\Jahr}[1]{#1} % Datum nur als Jahr
\newcommand{\Fachbegriff}[1]{\textbf{#1}}
\newcommand{\Klasse}[1]{\texttt{#1}}
\newcommand{\Methode}[1]{\texttt{#1}}
\newcommand{\MethodenRef}[1]{\texttt{#1}}
\newcommand{\Code}[1]{\textsf{#1}}

% text sizes
% \Huge
% \huge
% \LARGE
% \Large
% \large
% \normalsize (default)
% \small
% \footnotesize
% \scriptsize
% \tiny